%%%%%%%%%%%%%%%%%%%%%%%%%%%%%%%%%%%%%%%%%%%%%%%%%%%%%%%%%%%%%%%%%%%%%
%%                                                                 %%
%% Please do not use \input{...} to include other tex files.       %%
%% Submit your LaTeX manuscript as one .tex document.              %%
%%                                                                 %%
%% All additional figures and files should be attached             %%
%% separately and not embedded in the \TeX\ document itself.       %%
%%                                                                 %%
%%%%%%%%%%%%%%%%%%%%%%%%%%%%%%%%%%%%%%%%%%%%%%%%%%%%%%%%%%%%%%%%%%%%%

%%\documentclass[referee,sn-basic]{sn-jnl}% referee option is meant for double line spacing

%%=======================================================%%
%% to print line numbers in the margin use lineno option %%
%%=======================================================%%

%%\documentclass[lineno,sn-basic]{sn-jnl}% Basic Springer Nature Reference Style/Chemistry Reference Style

%%======================================================%%
%% to compile with pdflatex/xelatex use pdflatex option %%
%%======================================================%%

%%\documentclass[pdflatex,sn-basic]{sn-jnl}% Basic Springer Nature Reference Style/Chemistry Reference Style

%%\documentclass[sn-basic]{sn-jnl}% Basic Springer Nature Reference Style/Chemistry Reference Style
%%\documentclass[sn-mathphys]{sn-jnl}% Math and Physical Sciences Reference Style
%%\documentclass[sn-aps]{sn-jnl}% American Physical Society (APS) Reference Style
%%\documentclass[sn-vancouver]{sn-jnl}% Vancouver Reference Style
%%\documentclass[sn-apa]{sn-jnl}% APA Reference Style
%%\documentclass[sn-chicago]{sn-jnl}% Chicago-based Humanities Reference Style
%%\documentclass[sn-standardnature]{sn-jnl}% Standard Nature Portfolio Reference Style
\documentclass[default]{sn-jnl}% Default
% \documentclass[default,iicol]{sn-jnl}% Default with double column layout

%%%% Standard Packages
%%<additional latex packages if required can be included here>
%%%%
\usepackage{epsfig}
\usepackage{graphicx}
\usepackage{amsmath}
\usepackage{amssymb}
\usepackage{booktabs}
\usepackage{comment}
\usepackage{color}

\usepackage{latexsym}
\usepackage{makecell,diagbox}
\usepackage{algorithm}
\usepackage{listings}

\usepackage{amsfonts}
\usepackage{colortbl}
\usepackage{microtype}
\usepackage{multirow}
\usepackage{adjustbox}
\usepackage{enumitem}
\usepackage{float}
\usepackage{bm}
\usepackage{bbm} % \mathbbm

\usepackage{bbding}
\usepackage{pifont}
\usepackage{wasysym}
\usepackage{xcolor}
% \usepackage{caption}

%%%%%=============================================================================%%%%
%%%%  Remarks: This template is provided to aid authors with the preparation
%%%%  of original research articles intended for submission to journals published 
%%%%  by Springer Nature. The guidance has been prepared in partnership with 
%%%%  production teams to conform to Springer Nature technical requirements. 
%%%%  Editorial and presentation requirements differ among journal portfolios and 
%%%%  research disciplines. You may find sections in this template are irrelevant 
%%%%  to your work and are empowered to omit any such section if allowed by the 
%%%%  journal you intend to submit to. The submission guidelines and policies 
%%%%  of the journal take precedence. A detailed User Manual is available in the 
%%%%  template package for technical guidance.
%%%%%=============================================================================%%%%

\jyear{2022}%

%% as per the requirement new theorem styles can be included as shown below
\theoremstyle{thmstyleone}%
\newtheorem{theorem}{Theorem}%  meant for continuous numbers
%%\newtheorem{theorem}{Theorem}[section]% meant for sectionwise numbers
%% optional argument [theorem] produces theorem numbering sequence instead of independent numbers for Proposition
\newtheorem{proposition}[theorem]{Proposition}% 
%%\newtheorem{proposition}{Proposition}% to get separate numbers for theorem and proposition etc.

\theoremstyle{thmstyletwo}%
\newtheorem{example}{Example}%
\newtheorem{remark}{Remark}%

\theoremstyle{thmstylethree}%
\newtheorem{definition}{Definition}%

\raggedbottom
%%\unnumbered% uncomment this for unnumbered level heads

\begin{document}

\title[]{Response Letter}

%%=============================================================%%
%% Prefix	-> \pfx{Dr}
%% GivenName	-> \fnm{Joergen W.}
%% Particle	-> \spfx{van der} -> surname prefix
%% FamilyName	-> \sur{Ploeg}
%% Suffix	-> \sfx{IV}
%% NatureName	-> \tanm{Poet Laureate} -> Title after name
%% Degrees	-> \dgr{MSc, PhD}
%% \author*[1,2]{\pfx{Dr} \fnm{Joergen W.} \spfx{van der} \sur{Ploeg} \sfx{IV} \tanm{Poet Laureate} 
%%                 \dgr{MSc, PhD}}\email{iauthor@gmail.com}
%%=============================================================%%


\maketitle

%%%%%%%%% BODY TEXT

% We wish to extend our heartfelt appreciation to the editors and reviewers for their meticulous evaluation of our manuscript titled ``CLIP-Adapter: Better Vision-Language Models with Feature Adapters''. The insightful feedback are immensely valuable in enhancing the overall quality of our paper. 

We wish to extend our heartfelt appreciation to the editors and reviewers.
We have carefully incorporated all the suggestions and thoroughly revised our manuscript. The following section provides a concise summary of the modifications we have made and our detailed responses to address their concerns, point by point. 

\noindent\textcolor{red}{The revised paper is uploaded in ``Supplementary Material".}
%Accordingly, this document is structured as follows:
% \begin{itemize}
%     \item Section~\ref{sec:summary}: Summary of Changes

%     \item Section~\ref{sec:editor}: Response to Editor
    
%     % \item Section~\ref{sec:reviewer_1}: Response to Reviewer \#1
    
%     \item Section~\ref{sec:reviewer_3}: Response to Reviewer \#3
    
%     \item Section~\ref{sec:reviewer_4}: Response to Reviewer \#4
    
%     \item Annotated Revision (with highlighted text in red)
% \end{itemize}



% \section{Summary of Changes}
% \label{sec:summary}
% We carefully revised the paper to thoroughly address all questions and suggestions raised by the reviewers, providing point-by-point responses. Specifically, our revisions encompassed:

% \begin{enumerate}
%     \item In Section 4.2.5, we added more analysis about why adopting the visual adapter performs better than the text adapter for CLIP.
%     \item We provided more experimental details in the caption of Table 1 for a fair comparison.
%     \item We clarified the difference in image backbones between Table 1 and Table 2 to avoid confusion.
%     \item In Section 4.2.6, we added more reasons about why inserting adapters at the 12-th layer is better than all layers.
%     \item In Section 2, we added more references and discussions about four missing related works.
    
% \end{enumerate}


% \section{Response to Editor}
% \label{sec:editor}

% \textbf{Point 1:} The first-round revision has been reviewed by two experts in the field. There are no major issues; however several minor suggestions from reviewers still remain and are necessary to improve the quality of the manuscript. The authors are encouraged to prepare another round of minor revision to fix these minor comments.

% \noindent \textbf{Reply:} Thank you very much for your hard efforts on handling our submission. These constructive suggestions are highly helpful for us to improve our paper. We followed all the comments to revise our paper thoroughly. For further details, please refer to the following response.



\section{Response to Reviewer \#3}
\label{sec:reviewer_3}

\textbf{Q1:} Considering CoOp~\cite{zhou2021coop} as a method for fine-tuning the text encoder while maintaining a fixed image encoder, and viewing the ``CLIP-Adapter'' as an alternative approach, it is plausible to suggest that fine-tuning the image encoder of CLIP and keeping its text encoder fixed may yield superior results in transferring CLIP to downstream tasks. Do the authors possess additional insights regarding this matter?

\vspace{0.1cm}
\noindent\textbf{R1:} It is indeed that our CLIP-Adapter only with a visual adapter after CLIP's image encoder performs better than that with a text encoder or CoOp~\cite{zhou2021coop}, as compared in Table 1 and Figure 5 of the paper. However, we keep both the image and text encoders of CLIP frozen, and only train the lightweight visual or text adapters. As shown in Table 10 of the paper (Table~\ref{tab:finetune} below), fully fine-tuning CLIP's image encoder can achieve better performance than fine-tuning the text encoder, but still falls behind CLIP-Adapter that only trains the visual adapter, due to the over-fitting issue. Therefore, our approach not only attains better downstream accuracy, but also exhibits superior computational efficiency. Moreover, as we aim to tackle a non-trivial classification task for images, learning a visual adapter for CLIP's image encoder can exert more explicit and effective adaption performance than the textual branch.

\begin{table*}[h]
\centering
\caption{Comparison of finetuning different components of CLIP-Adapter on 16-shot ImageNet dataset with a single NVIDIA A100 GPU.}
\vspace*{5pt}
\begin{adjustbox}{width=0.7\linewidth}
	\begin{tabular}{ccccc}
	\toprule
		Visual Encoder &Textual Encoder
		 &Adapter &Accuracy & Train Time\\ \midrule
		- & - &\Checkmark &\textbf{61.33\%} &\textbf{50min}\\
		\Checkmark & - &\Checkmark &60.07\% &1h 20min\\
		- & \Checkmark &\Checkmark &57.88\% &3h+\\
	    \Checkmark & \Checkmark & \Checkmark & 52.78\% &3h+\\
	\bottomrule
	\end{tabular}
\end{adjustbox}
\label{tab:finetune}
\end{table*}


\section{Response to Reviewer \#4}
\label{sec:reviewer_4}

\noindent \textbf{Q1:} In Table 1, the authors compare the training time, parameters and classification accuracy with other methods. However, the authors should also provide the running dataset and training details, such as batch size and the number of samples, GPU memory and training time calculation are all related to these details. 

\vspace{0.1cm}
\noindent \textbf{R1:} Thanks for pointing this out! We conduct the comparison on ImageNet~\cite{deng2009imagenet} dataset with 16-shot training images, i.e., 16 samples per category. We adopt a batch size 32 for all compared methods on a single NVIDIA A100 GPU. We have added these details in the caption of Table 1.

% \vspace{5pt}
\vspace{0.4cm}

\noindent \textbf{Q2:} Table 2 shows that inserting CLIP-Adapter in layer 12 achieves a better tradeoff between training time and accuracy (70.88). However, Table 1 shows that the accuracy is about 60 for 16-shot ImageNet. The descriptions of ablations in Sec 4.2.6 are also unclear.

\vspace{0.1cm}
\noindent \textbf{R2:} Sorry for the confusion.
In Table 2, we conduct the ablation study with a ViT-B/16 image backbone, since only a transformer-based network can be inserted with adapters at different intermediate layers. However, in Table 1, we compare CLIP-Adapter with other existing methods with a ResNet-50 image backbone, which achieves lower results than ViT-B/16 (61.33\% vs. 70.88\%). We have clarified this point in the revised paper.


% \vspace{5pt}
\vspace{0.4cm}
\noindent \textbf{Q3:} Inserting CLIP-Adapter in all layers gets 67.67 accuracy. It would be better to give some reasons why the performance degrades.

\vspace{0.1cm}
\noindent \textbf{R3:} As shown in Table 2 of the paper (Table~\ref{tab:insert} below), inserting adapters in all layers yields a total of 5.20M parameters, which is much more weights than inserting only at the 12-th layer (0.52M). The latter can well alleviate the over-fitting issue on few-shot training data, and largely preserve the pre-trained knowledge of CLIP by inserting the adapter at the end.

\begin{table*}[h]
\centering
\caption{Comparison of adapters inserted into different layers of CLIP. We adopt ViT-B/16 as the image backbone for the visual adapter, and report the 16-shot ImageNet performance.}
\vspace*{5pt}
\begin{adjustbox}{width=0.8\linewidth}
	\begin{tabular}{lcccccccc}
	\toprule
		Insert Layers & All    & 12  &  10  &  8 & 6 & 4 & 2 &0 \\
		\midrule
		Accuracy (\%)  & 67.67		&70.88		&\textbf{71.85}	&70.55		&70.03		&70.14	   	&69.43		&69.05 \\
		GPU Memory (GiB)  & 5.98	 &	\textbf{2.22}		 &2.65	 &	2.89	 &	3.23		 &4.46		 &4.88		 &5.14 \\
            Parameters (M)  & 5.20	&	\textbf{0.52}		&0.78		&0.78	&	0.78		&0.78		&0.78		&0.78 \\
	\bottomrule
	\end{tabular}
\end{adjustbox}
\label{tab:insert}
\end{table*}


% \vspace{5pt}
\vspace{0.4cm}
\noindent \textbf{Q4:} Some references related to parameter-efficient fine-tuning are missing~\cite{lian2022scaling,hu2022sparse,sung2022lst,sun2022black}, which can be discussed in the revised version. There also exist amounts of typos in this version, which need to be addressed.

\vspace{0.1cm}
\noindent \textbf{R4:} Thanks for your suggestion. We have cited these excellent parameter-efficient fine-tuning works in the revised paper, and carefully discussed their relations with ours: ``Some other methods, e.g., black-box tuning~\cite{sun2022black}, ladder side-tuning~\cite{sung2022lst}, sparse structure search\cite{hu2022sparse}, and Scaling\&Shifting~\cite{lian2022scaling} also introduce different techniques for adapting large language and vision models in a parameter-efficient manner.
Different from existing approaches, the proposed CLIP-Adapter applies a simple residual transformation layer over the feature embedding or classifier weight generated by CLIP.''



%%===========================================================================================%%
%% If you are submitting to one of the Nature Portfolio journals, using the eJP submission   %%
%% system, please include the references within the manuscript file itself. You may do this  %%
%% by copying the reference list from your .bbl file, paste it into the main manuscript .tex %%
%% file, and delete the associated \verb+\bibliography+ commands.                            %%
%%===========================================================================================%%
\bibliographystyle{sn-basic}
\bibliography{custom}% common bib file
%% if required, the content of .bbl file can be included here once bbl is generated
%%\input sn-article.bbl

%% Default %%
%%\input sn-sample-bib.tex%

\end{document}
